\section{Relevamiento de datos necesarios}

Para diseñar la solución separamos los datos según el papel que tienen dentro del problema. No todos los datos tienen el mismo ciclo de vida ni se consultan de la misma forma. Por ejemplo, un pedido se consulta con filtros y relaciones exactas, mientras que un procedimiento escrito se recupera por similitud semántica.

\begin{table}[ht]
\centering
\begin{tabular}{p{0.24\linewidth}p{0.37\linewidth}p{0.29\linewidth}}
\toprule
Grupo & Datos principales & Uso dentro de la solución \\
\midrule
Operativos & Productos, categorías, clientes, segmentos, proveedores, pedidos, líneas de pedido, entregas e incidencias. & Responder consultas exactas del negocio y relacionar documentos con el contexto operativo. \\
Comerciales & Condiciones comerciales, precios, descuentos y períodos de vigencia. & Consultar qué condición corresponde a un cliente y producto en una fecha determinada. \\
Documentales & Documentos, versiones, clases documentales, sensibilidad, archivos y fragmentos. & Mantener el corpus, su procedencia, su vigencia y la unidad de recuperación. \\
Vectoriales & Modelos de embedding y vectores de cada fragmento. & Resolver búsquedas por similitud sobre documentación autorizada. \\
Acceso & Actores, perfiles autorizados y permisos por clase documental. & Limitar qué documentos puede recuperar cada perfil. \\
Interacción y auditoría & Consultas, respuestas, evidencias y eventos de auditoría. & Explicar respuestas anteriores y registrar acciones relevantes. \\
\bottomrule
\end{tabular}
\caption{Datos relevados para la solución.}
\label{tab:relevamiento}
\end{table}

Los datos operativos son necesarios porque el copiloto no debería resolver por similitud preguntas que tienen una respuesta exacta en la base. Por ejemplo, la condición comercial vigente se obtiene cruzando cliente, producto y fecha; los pedidos con entregas demoradas se obtienen relacionando pedidos, entregas e incidencias. Guardar estos datos de forma estructurada permite aplicar restricciones, realizar agregaciones y evitar ambigüedades.

Para el corpus documental necesitamos registrar más información que el texto. Cada documento tiene una clase, un nivel de sensibilidad, una procedencia y una o varias versiones. Cada versión conserva el archivo original mediante su ruta relativa, nombre, tipo MIME, tamaño y checksum SHA-256. De esta manera, la base no guarda el archivo completo, pero puede verificar que el archivo al que hace referencia es el esperado.

Cada versión se divide en fragmentos. El fragmento es la unidad que se vectoriza y también la unidad que luego se puede mostrar como evidencia. Esto permite indicar con precisión qué sección de qué versión fue utilizada para fundamentar una respuesta, en lugar de asociar la respuesta a un documento completo.

Los datos de acceso son importantes porque no todos los perfiles pueden ver las mismas clases documentales. El permiso se modela como una relación entre perfil y clase documental. La ausencia de un permiso significa denegación, por lo que no se necesitan listas de excepciones individuales. Finalmente, los datos de interacción y auditoría permiten reconstruir quién consultó, con qué perfil, qué resultado obtuvo y qué evidencia se utilizó, sin duplicar el contenido sensible de los documentos en el registro de auditoría.

Para la implementación usamos datos sintéticos. Esto evita trabajar con datos personales, comerciales o documentos internos reales y permite repetir la carga con los mismos conteos, archivos, vectores y resultados esperados.
